\begin{definition}
Let $\mathbb G$ be a sub-elliptic setting (\noderef{CarnotSetting}) with group law $*$ and
homogeneous norm $\|\cdot\|$, and write $d(x,y):=\|y^{-1}*x\|$ for the left-invariant
distance of Section~2.4. For $\Om\subseteq\mathbb R^{n}$ and $\delta\in\mathbb R$ the
\emph{inner domain} $\Om_\delta$ of (3.8) is
\[
 \Om_\delta:=\big\{x\in\Om:\ \delta<d(x,y)\ \text{for every}\ y\in\del\Om\big\} .
\]
The paper writes this as $\{x\in\Om:\ \dist(x,\del\Om)>\delta\}$ with
$\dist(x,\del\Om)=\inf\{d(x,y):y\in\del\Om\}$, and the two agree: whenever $\Om$ is
bounded, $\del\Om$ is compact, so for $\del\Om\ne\emptyset$ the infimum is attained and
"$\dist(x,\del\Om)>\delta$" is equivalent to "$\delta<d(x,y)$ for every $y\in\del\Om$";
and for $\del\Om=\emptyset$ the displayed condition is vacuous, giving $\Om_\delta=\Om$,
which matches the paper's convention $\inf\emptyset=+\infty$. The quantified form is used
here because it is the form in which the condition is consumed: every use below reads it as
"every boundary point is $d$-far from $x$".
\end{definition}
